\chapter{MỞ ĐẦU}
\label{ch:mo-dau}

\section{Lý do chọn đề tài}
\label{sec:ly-do}

Thị trường chứng khoán là một hệ thống chịu ảnh hưởng đồng thời của thông tin kinh tế,
hoạt động của doanh nghiệp, kỳ vọng của nhà đầu tư và trạng thái chung của thị trường.
Trong số các nguồn thông tin đó, tin tức tài chính có thể làm thay đổi kỳ vọng trong
thời gian ngắn, đặc biệt khi tin đề cập đến kết quả kinh doanh, thay đổi chính sách,
hoạt động huy động vốn hoặc các sự kiện bất thường của doanh nghiệp. Vì vậy, việc
chuyển một dòng tin văn bản thành tín hiệu định lượng và kiểm tra xem tín hiệu đó có
bổ sung thông tin cho chuỗi giá hay không là một bài toán có ý nghĩa cả về xử lý ngôn
ngữ tự nhiên lẫn phân tích chuỗi thời gian. Các nghiên cứu quốc tế đã khai thác cả điểm
cảm xúc và biểu diễn ngữ cảnh của tin tức để dự báo chuyển động cổ phiếu, nhưng kết quả
phụ thuộc mạnh vào thị trường, ngôn ngữ, cách gán nhãn và giao thức đánh giá của từng
nghiên cứu \cite{Halder2022_221107392v1,Chen2021_210708721v1}.

\subsection{Bối cảnh tiếng Việt và thị trường Việt Nam}

Đối với thị trường Việt Nam, phần lớn tin tức tài chính được công bố bằng tiếng Việt.
Tiếng Việt có đặc trưng về từ ghép, tách từ, dấu thanh, cách viết tên doanh nghiệp và
cách biểu đạt con số tài chính. Những đặc trưng này làm cho việc sử dụng trực tiếp một
mô hình được huấn luyện chủ yếu trên tiếng Anh không phải là lựa chọn hiển nhiên. PhoBERT
được xây dựng như một mô hình ngôn ngữ tiền huấn luyện cho tiếng Việt, do đó cung cấp
một backbone phù hợp hơn cho các bài toán phân loại văn bản trong miền ngôn ngữ này
\cite{Nguyen2020_200300744v3}. Ở miền tiếng Anh, FinBERT cho thấy việc thích nghi
backbone BERT với ngữ liệu tài chính có thể tạo ra biểu diễn hữu ích cho phân tích cảm
xúc tài chính \cite{Yang2020_200608097v2}. Hai hướng phát triển này gợi ý một nguyên
tắc thiết kế cho đề tài: sử dụng mô hình ngôn ngữ phù hợp với ngôn ngữ của dữ liệu,
đồng thời kiểm tra riêng chất lượng phân loại cảm xúc trước khi đưa tín hiệu đó vào
mô hình dự báo giá.

Thị trường HOSE và nhóm cổ phiếu VN30 được chọn làm bối cảnh nghiên cứu vì đây là một
phạm vi có lịch sử giao dịch tương đối dày, các mã có mức độ quan tâm thông tin cao và
phù hợp với việc xây dựng dữ liệu theo ngày. Tuy nhiên, phạm vi này không đại diện cho
toàn bộ thị trường chứng khoán Việt Nam. Mỗi mã có đặc điểm ngành, thanh khoản và mật
độ tin tức khác nhau, còn quan hệ giữa tin và giá có thể thay đổi theo từng giai đoạn.
Do đó, kết quả của đề tài được diễn giải như bằng chứng thực nghiệm trong phạm vi các
mã được chọn, không phải như một khẳng định tổng quát cho mọi cổ phiếu.

\subsection{Động cơ kết hợp hai loại dữ liệu}

Dữ liệu giá chứa thông tin về phản ứng đã được thị trường thể hiện qua các đại lượng
như giá mở cửa, giá cao nhất, giá thấp nhất, giá đóng cửa, khối lượng và lợi suất.
Ngược lại, tin tức có thể chứa tín hiệu về sự kiện hoặc kỳ vọng trước khi phản ứng đó
được phản ánh đầy đủ trong chuỗi giá. Hai nguồn dữ liệu vì vậy có tính bổ sung, nhưng
không đồng nhất về dạng biểu diễn, thời điểm xuất hiện và độ tin cậy. Một mô hình chỉ
dùng giá có thể bỏ qua thông tin văn bản, trong khi một mô hình chỉ dùng cảm xúc có thể
bỏ qua xu hướng, biến động và trạng thái kỹ thuật của cổ phiếu.

Các kiến trúc kết hợp tin tức và giá trong nghiên cứu quốc tế thường đi từ việc trích
xuất sentiment bằng FinBERT rồi đưa vào LSTM, đến việc sử dụng embedding ngữ cảnh hoặc
LLM để mô hình hóa thông tin văn bản giàu hơn \cite{Halder2022_221107392v1,
Gu2024_240716150v1,Elahi2024_241101368v1}. Tuy vậy, một mức cải thiện được báo cáo
trên dữ liệu nước ngoài không tự động chuyển thành hiệu quả trên HOSE. Đề tài vì thế
đặt trọng tâm vào phép so sánh có kiểm soát giữa mô hình chỉ dùng giá và mô hình có
thêm sentiment, dùng cùng cửa sổ thời gian, cùng cách tiền xử lý và cùng tập kiểm tra.

\subsection{Vấn đề phương pháp cần giải quyết}

Khó khăn cốt lõi không chỉ nằm ở việc chọn một mô hình sâu hơn. Có ít nhất bốn vấn đề
phương pháp cần được xử lý đồng thời.

Thứ nhất, văn bản phải được gắn với đúng mã cổ phiếu và đúng thời điểm. Một bài viết
được cập nhật sau thời điểm dự báo không thể được dùng để giải thích quyết định của
phiên trước đó. Nếu tin sau giờ giao dịch, tin cuối tuần hoặc tin có timestamp không
được xử lý nhất quán, mô hình có thể đạt điểm số cao do nhìn thấy thông tin tương lai.

Thứ hai, tín hiệu cảm xúc ở cấp bài viết phải được tổng hợp thành đặc trưng ở cấp mã và
ngày. Việc chỉ lấy một bài nổi bật hoặc ghép toàn bộ tin trong ngày mà không xét số
lượng, tỷ lệ lớp và ngày không có tin có thể làm mất thông tin về cường độ và độ phủ
của tín hiệu.

Thứ ba, dữ liệu tài chính thường mất cân bằng và không dừng. Một mô hình có thể đạt
accuracy tương đối cao chỉ nhờ dự đoán lớp phổ biến, trong khi bỏ sót lớp giảm hoặc
lớp đi ngang. Đây là lý do đề tài dùng macro-F1 và balanced accuracy bên cạnh các chỉ
số theo lớp.

Thứ tư, phép đánh giá phải tôn trọng thứ tự thời gian. Chuẩn hóa, điền khuyết, chọn
đặc trưng, chọn ngưỡng nhãn hoặc tinh chỉnh mô hình bằng dữ liệu tương lai đều làm suy
giảm giá trị của kết quả ngoài mẫu. Các công trình về dự báo chuyển động cổ phiếu cũng
cho thấy việc thiết kế benchmark, dữ liệu và cách đánh giá có vai trò quan trọng không
kém việc thay đổi kiến trúc \cite{Feng2018_181009936v2,Chen2021_210708721v1}.

\subsection{Ý nghĩa của đề tài}

Ý nghĩa của đề tài được xác định theo hai hướng. Về học thuật, đề tài xây dựng một
quy trình nối từ mô hình ngôn ngữ tiếng Việt đến mô hình chuỗi thời gian, qua đó kiểm
tra một cách định lượng xem sentiment có tạo ra phần thông tin gia tăng khi dự báo
xu hướng phiên kế tiếp hay không. Về kỹ thuật, đề tài tạo ra một pipeline có thể tái
lập gồm thu thập dữ liệu, kiểm tra chất lượng, xử lý point-in-time, trích xuất xác suất
cảm xúc, tổng hợp theo mã và ngày, huấn luyện mô hình, đánh giá walk-forward và ghi
nhận cấu hình.

Đề tài không đặt mục tiêu chứng minh rằng sentiment luôn làm mô hình tốt hơn. Trong
dữ liệu tài chính, tín hiệu văn bản có thể yếu, trễ, trùng với phản ứng giá hoặc chỉ có
ích trong một số giai đoạn. Một kết quả không cải thiện, hoặc cải thiện nhỏ và không
ổn định, vẫn có giá trị nếu được thu nhận bằng giao thức hợp lệ và được trình bày cùng
giới hạn. Cách đặt vấn đề này phù hợp với hướng nghiên cứu gần đây, trong đó giá trị
gia tăng của sentiment được xem là một giả thuyết cần kiểm tra thay vì một kết luận
được giả định trước \cite{Jiang2023_230602136v3,Siala2026_260200086v3}.

\section{Mục tiêu nghiên cứu}
\label{sec:muc-tieu}

Mục tiêu tổng quát của đồ án là xây dựng và đánh giá một hệ thống dự báo xu hướng biến
động giá cổ phiếu cho thị trường HOSE. Hệ thống sử dụng PhoBERT để tạo tín hiệu cảm xúc
từ tin tức tài chính tiếng Việt, sử dụng đặc trưng kỹ thuật để biểu diễn chuỗi giá,
sau đó hợp nhất hai nguồn thông tin trong mô hình dự báo ba lớp cho phiên kế tiếp.
Trọng tâm của nghiên cứu là đo lường phần đóng góp biên của sentiment dưới giao thức
đánh giá chống rò rỉ, không phải tối ưu một con số hiệu năng duy nhất.

\subsection{Câu hỏi nghiên cứu}

Đề tài trả lời ba câu hỏi nghiên cứu sau:

\begin{enumerate}[nosep]
  \item Mô hình PhoBERT được tinh chỉnh trên dữ liệu tin tức tài chính tiếng Việt có
    thể phân biệt ba lớp tích cực, trung tính và tiêu cực ở mức nào? Chất lượng có thay
    đổi ra sao khi so sánh head tuyến tính với một biến thể tăng cường có cùng dữ liệu?
  \item Việc bổ sung đặc trưng cảm xúc vào mô hình dự báo có cải thiện macro-F1,
    balanced accuracy và macro OvR-AUC so với mô hình chỉ sử dụng giá hay không? Nếu có,
    mức chênh lệch có ổn định qua các cửa sổ thời gian hay chỉ xuất hiện ở một số fold?
  \item Tín hiệu cảm xúc có biểu hiện khác nhau giữa các giai đoạn thị trường, giữa
    ngày có nhiều tin và ngày ít tin, hoặc giữa ngày gắn với sự kiện và ngày thông thường
    hay không?
\end{enumerate}

\subsection{Mục tiêu cụ thể}

Để trả lời các câu hỏi trên, đồ án đặt ra các mục tiêu có thể kiểm chứng:

\begin{enumerate}[nosep]
  \item Xây dựng bộ dữ liệu tin tức tài chính tiếng Việt và dữ liệu giá điều chỉnh của
    một rổ 10 mã thuộc VN30 trong giai đoạn từ năm 2020 đến quý I năm 2026, kèm nguồn,
    timestamp, mã cổ phiếu và thông tin provenance.
  \item Thiết kế quy trình làm sạch, loại trùng, ánh xạ bài viết theo mã và phiên giao
    dịch, trong đó dữ liệu chỉ được sử dụng nếu phù hợp với thời điểm dự báo.
  \item Tinh chỉnh PhoBERT cho bài toán phân loại ba lớp, sinh xác suất của từng lớp
    cho mỗi bài viết và tổng hợp thành vector sentiment theo mã và ngày.
  \item Trích xuất lợi suất và các chỉ báo kỹ thuật từ OHLCV điều chỉnh, chuẩn hóa các
    đặc trưng bằng tham số chỉ học từ phần huấn luyện của từng cửa sổ.
  \item Xây dựng thang mô hình gồm majority, random, mô hình cổ điển chỉ dùng giá,
    LSTM chỉ dùng giá và LSTM hai nhánh có fusion giữa giá với sentiment. TFT được
    xem là tầng mở rộng sau khi graduation floor đã được kiểm tra.
  \item Đánh giá dự báo bằng walk-forward có holdout cuối được khóa, báo cáo macro-F1
    là chỉ số chính, balanced accuracy, macro OvR-AUC và các chỉ số theo lớp là chỉ số
    bổ sung, đồng thời ước lượng khoảng tin cậy cho chênh lệch ablation.
  \item Ghi nhận cấu hình, seed, phiên bản thư viện, phiên bản dữ liệu và số lần thử,
    để kết quả có thể được tái lập hoặc kiểm tra lại khi có thay đổi dữ liệu.
\end{enumerate}

\subsection{Nguyên tắc diễn giải kết quả}

Các mốc hiệu năng được nêu trong tài liệu nguồn, chẳng hạn kết quả phân loại của một
seed dữ liệu hoặc mức directional accuracy của một nghiên cứu nước ngoài, chỉ được
xem là tham chiếu kỹ thuật. Chúng được đo trên ngôn ngữ, thị trường, khoảng thời gian
và cách chia dữ liệu khác với đề tài, nên không được dùng như ngưỡng đạt hoặc lời hứa
cho HOSE. Tương tự, attention hoặc variable importance chỉ mô tả những tín hiệu mà
mô hình đang sử dụng; chúng không đủ để kết luận sentiment là nguyên nhân gây ra biến
động giá.

Vì mục tiêu là kiểm tra giá trị gia tăng, mọi phát biểu về việc sentiment hữu ích hay
không hữu ích phải dựa trên các cặp mô hình được huấn luyện và kiểm tra dưới cùng một
protocol. Một kết quả thuyết phục cần nêu rõ chênh lệch điểm số, độ biến thiên giữa các
fold, khoảng tin cậy, phân bố lớp, phạm vi mã và các giới hạn của dữ liệu.

\section{Đối tượng và phạm vi nghiên cứu}
\label{sec:pham-vi}

\subsection{Đối tượng nghiên cứu}

Đối tượng trung tâm của đề tài là mối quan hệ dự báo giữa cảm xúc của tin tức tài chính
tiếng Việt và xu hướng lợi suất của cổ phiếu trong phiên kế tiếp. Mối quan hệ này được
nghiên cứu thông qua hai loại dữ liệu. Loại thứ nhất là văn bản tin tức, gồm tiêu đề,
nội dung khi có thể thu thập, thời điểm đăng hoặc cập nhật, nguồn và mã cổ phiếu được
ánh xạ. Loại thứ hai là dữ liệu thị trường OHLCV theo ngày, trong đó giá được xử lý theo
quy ước điều chỉnh để hạn chế ảnh hưởng của các sự kiện doanh nghiệp.

Về mô hình, đối tượng nghiên cứu gồm một backbone NLP tiếng Việt là PhoBERT và các
biến thể head dùng để phân loại cảm xúc, cùng nhóm mô hình chuỗi thời gian gồm mô hình
cổ điển, LSTM và Temporal Fusion Transformer. PhoBERT được chọn vì được tiền huấn luyện
cho tiếng Việt, còn việc tham khảo FinBERT giúp đặt lựa chọn này trong dòng nghiên cứu
NLP tài chính quốc tế \cite{Nguyen2020_200300744v3,Yang2020_200608097v2}.

\subsection{Phạm vi thị trường và mã cổ phiếu}

Phạm vi thị trường là sàn HOSE. Rổ nghiên cứu gồm 10 mã thuộc VN30, được lựa chọn theo
tiêu chí thanh khoản, độ phủ tin và khả năng thu thập chuỗi giá liên tục. Danh sách mã
và ngày bắt đầu sử dụng phải được lưu trong manifest, vì thành phần VN30 có thể thay
đổi theo thời gian. Trong các bước chia dữ liệu, các mã được giữ trong cùng lịch phiên
để việc so sánh theo ngày không bị lệch giữa các mã.

Việc chọn 10 mã là một giới hạn có chủ đích. Phạm vi vừa đủ để kiểm tra pipeline đa mã
và giảm phụ thuộc vào một case study duy nhất, nhưng chưa đủ để đại diện cho mọi nhóm
ngành hoặc mọi mức thanh khoản trên HOSE. Do đó, kết quả cần được gắn với rổ mã và
khung thời gian cụ thể thay vì diễn giải như kết luận cho toàn thị trường.

\subsection{Phạm vi thời gian và tần suất}

Dữ liệu được giới hạn từ năm 2020 đến hết quý I năm 2026, tùy ngày cuối có dữ liệu hợp
lệ của từng nguồn. Đơn vị quan sát chính là ngày giao dịch. Mục tiêu tại ngày $t$ là
dự báo lớp của lợi suất từ ngày $t$ sang phiên giao dịch kế tiếp. Tần suất ngày phù hợp
với mật độ tin tức và tránh mở rộng bài toán sang các vấn đề vi cấu trúc của dữ liệu
trong ngày.

Mốc cutoff chính của protocol là 15 giờ theo múi giờ Việt Nam. Tin sau cutoff, tin vào
cuối tuần hoặc ngày nghỉ được gán sang phiên giao dịch kế tiếp theo một quy tắc bảo
thủ. Các trường hợp timestamp thiếu hoặc mâu thuẫn không được tự động suy đoán mà phải
được đánh dấu, loại khỏi tập phân tích chính hoặc báo cáo riêng.

\subsection{Phạm vi bài toán và giới hạn}

Bài toán dự báo là phân loại ba lớp:

\begin{itemize}[nosep]
  \item \textbf{UP}: lợi suất phiên kế tiếp lớn hơn ngưỡng tăng;
  \item \textbf{FLAT}: lợi suất nằm trong vùng trung tính;
  \item \textbf{DOWN}: lợi suất nhỏ hơn ngưỡng giảm.
\end{itemize}

Ngưỡng được ước lượng trong tập huấn luyện của từng cửa sổ, không dùng phân vị của
toàn bộ giai đoạn. Quy ước này giúp tránh việc phân phối lợi suất của tương lai ảnh
hưởng đến nhãn của quá khứ. Đề tài không dự báo giá tuyệt đối, không nghiên cứu giao
dịch tần suất cao, sổ lệnh, phái sinh hoặc khuyến nghị đầu tư. Backtesting chiến lược,
Sharpe và maximum drawdown được dành cho hướng phát triển, không phải tiêu chí kết luận
của graduation floor.

\subsection{Đóng góp dự kiến}

Trong phạm vi được xác định, các đóng góp dự kiến gồm:

\begin{enumerate}[nosep]
  \item Một pipeline thực nghiệm tiếng Việt nối dữ liệu tin tức với dữ liệu giá theo
    mã và thời điểm, có kiểm tra deduplication, provenance và cutoff.
  \item Một quy trình dùng PhoBERT để tạo phân phối xác suất cảm xúc thay vì chỉ dùng
    một nhãn rời rạc, sau đó tổng hợp tín hiệu theo mã và ngày.
  \item Một kiến trúc sàn gồm LSTM hai nhánh và fusion bằng concat, đặt cạnh các
    baseline chỉ giá để đo phần đóng góp biên của sentiment.
  \item Một giao thức walk-forward có train-only preprocessing, holdout khóa cuối,
    macro-F1 và bootstrap CI, giúp tách kết quả mô hình khỏi các lỗi rò rỉ dữ liệu.
\end{enumerate}

Các đóng góp trên là đóng góp về hệ thống thực nghiệm và khả năng kiểm chứng trong
phạm vi dữ liệu đã chọn. Đồ án không tuyên bố tạo ra một mô hình tối ưu cho mọi thị
trường, không tuyên bố quan hệ nhân quả giữa tin và giá, và không xem điểm số của các
công trình trên dữ liệu nước ngoài là bằng chứng trực tiếp cho dữ liệu Việt Nam.

\section{Cấu trúc đồ án}
\label{sec:cau-truc}

Phần còn lại của đồ án được tổ chức theo trình tự từ cơ sở vấn đề, khảo sát nghiên cứu,
thiết kế phương pháp đến đánh giá và giới hạn.

\subsection{Chương 2, Tổng quan}

Chương~\ref{ch:tong-quan} phân tích ba trụ nghiên cứu. Trụ thứ nhất là sentiment tài
chính và backbone ngôn ngữ, từ FinBERT trong tiếng Anh đến PhoBERT trong tiếng Việt.
Trụ thứ hai là mô hình chuỗi thời gian và Transformer cho dự báo chuyển động. Trụ thứ
ba là hợp nhất tin tức với giá. Cuối chương, các hướng được so sánh theo dữ liệu, cách
biểu diễn, mô hình, cách đánh giá và khả năng tái lập, từ đó xác định khoảng trống mà
đề tài tập trung xử lý.

\subsection{Chương 3, Phương pháp}

Chương~\ref{ch:phuong-phap} trình bày cơ sở lý thuyết của Transformer, self-attention,
PhoBERT, LSTM và TFT; sau đó mô tả kiến trúc hai nhánh, nguồn dữ liệu, tiền xử lý,
gán nhãn, tạo đặc trưng, thang mô hình và giao thức đánh giá chống rò rỉ. Các quy tắc
cutoff, chia dữ liệu và fit tham số được viết tường minh để người đọc có thể kiểm tra
điểm nào của pipeline có nguy cơ nhìn thấy tương lai.

\subsection{Chương 4, Kết quả thực nghiệm}

Chương~\ref{ch:ket-qua} mô tả thiết lập thực nghiệm, thống kê dữ liệu, chất lượng mô
hình sentiment, kết quả dự báo và bảng ablation. Các ô kết quả chưa được đo được giữ
ở dạng trống để điền từ manifest thực nghiệm. Phần bàn luận được tổ chức quanh ba câu
hỏi nghiên cứu, với nguyên tắc không thay thế số liệu thực bằng mốc tham chiếu từ
công trình khác.

\subsection{Chương 5 và Chương 6}

Chương~\ref{ch:ket-luan} tổng hợp những gì có thể kết luận trong phạm vi protocol,
phân biệt kết quả đã đo với mục tiêu thiết kế và nêu các giới hạn. Chương~\ref{ch:huong-phat-trien}
đề xuất các hướng mở rộng theo thứ tự ưu tiên, gồm đánh giá giá trị kinh tế, mở rộng
ngữ liệu và mã, cải tiến fusion, xử lý concept drift và tăng khả năng giải thích.
